\begin{table}[!htbp]
\centering
\small
\renewcommand{\arraystretch}{1.25}
\begin{tabularx}{\textwidth}{>{\raggedright\arraybackslash}p{2.6cm}>{\raggedright\arraybackslash}p{5.5cm}>{\raggedright\arraybackslash}X}
\toprule
\textbf{Giai đoạn vòng đời} & \textbf{Lệnh OpenSSL thực thi trong Lab} & \textbf{Ca kiểm thử liên quan} \\
\midrule
\textbf{1. Sinh cặp khóa} & \texttt{openssl genrsa -aes256 ...} hoặc \texttt{openssl ecparam ...} & Nền tảng mọi kịch bản \\
\midrule
\textbf{2. Tạo CSR (PoP)} & \texttt{openssl req -new -key ... -out ...} & Tạo yêu cầu ký chứng thư kèm PoP \\
\midrule
\textbf{4. Cấp phát chứng thư} & \texttt{openssl ca -in ... -out ... -extfile ...} & Cấp phát X.509 với extension; TC05 (chữ ký bị sửa đổi bit), TC08 (cấp sai cờ EKU) \\
\midrule
\textbf{5. Triển khai cấu hình} & Cấu hình tệp chứng thư và chuỗi Intermediate CA bundle & TC06 (thiếu Intermediate CA bundle) \\
\midrule
\textbf{6. Xác minh đường dẫn} & \texttt{openssl verify -CAfile ... -untrusted ... -purpose ... -verify\_hostname ...} & TC01 (thẩm định thành công toàn diện), TC02 (self-signed không thuộc trust store), TC04 (hostname không khớp SAN), TC08 (sai mục đích EKU), TC10 (CA ngoài trust store) \\
\midrule
\textbf{8. Kiểm tra thời hạn} & \texttt{openssl verify -attime ...} (hoặc so khớp thời gian thực) & TC03 (thời điểm kiểm tra nằm ngoài hạn hiệu lực) \\
\midrule
\textbf{9. Thu hồi chứng thư} & \texttt{openssl ca -revoke ...}, \texttt{openssl ca -gencrl ...} và xác minh với \texttt{-crl\_check} & TC07 (chứng thư bị thu hồi đưa vào CRL) \\
\bottomrule
\end{tabularx}
\caption{Ánh xạ các giai đoạn vòng đời chứng thư X.509 với lệnh OpenSSL và ca kiểm thử.}
\label{tab:lifecycle-openssl-tc}
\end{table}
